\section{Formal scope of primitive content and power extraction}
\label{sec-collective-formal-scope}

The ordinary arguments in the preceding new sections were completed before
their selected integer and real-arithmetic consequences were formalized.
The present formal supplement contains forty new declarations:
twenty-four in \texttt{ComplementContent.lean} and sixteen in
\texttt{SquarefreePowerExtraction.lean}.
All forty were freshly compiled together with fifty-five unchanged
dependency declarations. Every one of these ninety-five declarations has
an explicit axiom query. The axiom union is precisely
\texttt{propext}, \texttt{Classical.choice}, and \texttt{Quot.sound}.
The compiler is Lean 4.32.0 and the Mathlib revision is
\texttt{81a5d257c8e410db227a6665ed08f64fea08e997}.
Only the Mathlib compiled cache is reused; the six listed project modules
are compiled from their recorded source bytes in a fresh directory.

The first module uses the actual integer-pair Eisenstein multiplication,
its positive norm, the integer gcd of the coordinates, and exact integer
division to define primitive reduction. It constructs the actual
normalizer \(c d\) and proves its divisibility into \(N(V)\) by B\'ezout.
It then proves the three recovered phase coordinates, their exact common
gcd, and equality of the primitive absolute boundary product. Its radical
is the product of the actual natural prime factors; for \(T>0\), its signed ledger is
\[
 \sum_{p\mid T}\bigl(v_p(T)-3\bigr)\log p.
\]
Both functions are therefore exactly unchanged. The cleared phase
equalities and the integrality of the common multiplier are proved from
the actual map, not supplied as hypotheses to its final theorem.

The second module starts with the explicit natural-number profile
\[
 RD=VQ^g,\quad R\text{ squarefree},\quad (R,D)=1,\quad g\ge2.
\]
It proves \(R\mid V\) and \(Q^g\mid D\) without requiring
\((V,Q)=1\). For \(Q>1\), an actual prime divisor of \(Q\) bounds
\(g\) by its multiplicity in \(D\). The module derives both logarithmic
compression inequalities and, from the explicitly stated prime profile
and mass hypotheses, the lower bound
\[
 \frac{\log V}{g}\ge\frac{\delta\log7}{6}B.
\]
The hypotheses realizing this profile in the interval product, and its
large-\(B\) mass estimate, remain the preceding ordinary proofs. They
have not been made formal by passing their finite consequences.

Eight finite replay records cover the collective content, complement
determinants, primitive power profile, elliptic local-entry coordinates,
actual divisor-partition witnesses, the recursive Lucas certificates for
the large-prime partition example, the height-transport check, and the
affine selectors' exact depths, content, finite local roots and kernels.
The portable height check recomputes the exact rational leading-digit
calculations and checks algebra on the recorded precision balls; it does
not recompute those balls with PARI. The separate full PARI 2.15.4 replay
was actually run locally and is distinguished from this portable check.
These replays do not prove prime
asymptotics, an infinite family, a formal-logarithm theorem, Jacobian
closure, a height pairing, or rational-point completeness. Those results
retain their ordinary proofs and source attributions. The independent
reviews are internal mathematical reviews, not external peer review.

The completed statements do not bound the full boundary signed tail,
cover every actual root or every ABC triple, compute a complete quadratic
Chabauty set, or establish uniform point heights as exponents and residuals
vary. Standard ABC remains unproved and undisproved. In particular, the
input-norm height gain cannot be counted as boundary radical credit;
the exact complement theorem identifies one specific construction where
that proposed credit disappears under primitive normalization.
